\paragraph{有理切片障碍、极化缺失与 \(F_8/F_9\) 隐通道分离}
在前述推前谱满对称性（定理~\ref{thm:terminal-window6-pushforward-charpoly-galois}、推论~\ref{cor:terminal-window6-pushforward-single-galois-orbit}）、有理交换子幂等刚性（定理~\ref{thm:terminal-windowm-rational-commutant-idempotent-rigidity}）、masa 交换子 \(*\)-代数（定理~\ref{thm:terminal-window6-pushforward-commutant-masa}）、强可并性失败的 \(F_8\) 二值见证（命题~\ref{prop:terminal-foldbin6-strong-lumpability-fails} 与推论~\ref{cor:terminal-foldbin6-strong-lumpability-binary-witness}），以及 boundary 扇区纯 \(F_9\) 二层别名（推论~\ref{cor:terminal-foldbin6-boundary-pure-f9-alias}）之外，window-\(6\) 还呈现一条更细的动力学分层：非平凡慢模切片不能在 \(\QQ\)-代数层实现，centered visible sector 不支持任何 \(\mathsf T_6\)-等变的复极化或辛极化，而最小的一阶 Markov 修复又完全落在被截断的 \(F_8\) 尾位通道，并且与 canonical boundary \(F_9\)-sheet parity 严格分离。

\begin{corollary}[window-\(6\) 的非平凡谱切片不存在有理谱投影]\label{cor:window6-no-rational-spectral-slice-projector}
记
\[
\Lambda_6:=\mathrm{spec}(P_6)\setminus\{1\},
\]
并对任意非空真子集 \(S\subsetneq \Lambda_6\)，记 \(E_S\) 为复化空间上投影到
\[
\bigoplus_{\lambda\in S}\ker(P_6-\lambda I)
\]
的谱投影。则
\[
E_S\notin M_{21}(\QQ).
\]
更强地，不存在任何 \(f\in\QQ[x]\) 使
\[
f(\lambda)=1\quad(\lambda\in S),\qquad
f(\lambda)=0\quad(\lambda\in \Lambda_6\setminus S).
\]
\end{corollary}

\begin{proof}
若 \(E_S\in M_{21}(\QQ)\)，则 \(E_S^2=E_S\) 且 \(E_S P_6=P_6 E_S\)。
由推论~\ref{cor:terminal-window6-rational-observable-compression-triviality}，与 \(P_6\) 交换的有理幂等元仅可能为
\[
0,\qquad I,\qquad \Pi_6,\qquad I-\Pi_6,
\]
其中 \(\Pi_6\) 为稳态方向上的投影。
然而当 \(S\subsetneq\Lambda_6\) 为非空真子集时，\(E_S\) 既不为 \(0\) 或 \(I\)，也不对应于仅选取平凡本征值 \(1\) 或其补空间的投影，故矛盾。

若存在 \(f\in\QQ[x]\) 满足所述 \(0/1\) 取值条件，则 \(f(P_6)\) 恰为 \(E_S\)；这将产生一个与 \(P_6\) 交换的有理谱投影，已被前半部分排除。证毕。
\end{proof}

\leanverified{paper\_window6\_no\_rational\_spectral\_slice\_projector}
\leanverified{paper\_window6\_no\_equivariant\_complex\_or\_symplectic\_polarization}
\begin{theorem}[centered visible sector 上不存在 \texorpdfstring{$\mathsf T_6$}{T6}-等变复极化或辛极化]\label{thm:window6-no-equivariant-complex-or-symplectic-polarization}
令
\[
H_6:=\ell^2(X_6,\pi_6),\qquad
H_6^0:=\{f\in H_6:\langle f,\mathbf 1\rangle_{\pi_6}=0\},
\]
并把 \(H_6^0\) 视为实内积空间。则：
\begin{enumerate}
  \item 不存在实线性算子 \(J:H_6^0\to H_6^0\) 使
  \[
  J^2=-I,\qquad J\mathsf T_6=\mathsf T_6J.
  \]
  \item 不存在非退化交替双线性型 \(\Omega\) 于 \(H_6^0\) 上，使
  \[
  \Omega(\mathsf T_6u,v)=\Omega(u,\mathsf T_6v)\qquad(\forall\,u,v\in H_6^0).
  \]
\end{enumerate}
\leanverified{paper\_window6\_no\_equivariant\_complex\_or\_symplectic\_polarization}
\end{theorem}

\begin{proof}
由定理~\ref{thm:terminal-window6-pushforward-commutant-masa}，\(\mathsf T_6\) 具有简单实谱，从而
\[
H_6^0=\bigoplus_{\lambda\in\Lambda_6}E_\lambda,
\]
其中每个 \(E_\lambda=\ker(\mathsf T_6-\lambda I)\) 都是一维实特征子空间。

若 \(J\mathsf T_6=\mathsf T_6J\)，则 \(J\) 保持每个 \(E_\lambda\)。
由于 \(E_\lambda\) 一维，必有 \(J|_{E_\lambda}=a_\lambda I\) 对某个 \(a_\lambda\in\RR\)。
于是
\[
J^2|_{E_\lambda}=a_\lambda^2 I,
\]
不可能等于 \(-I\)。故第 \(1\) 条成立。

若 \(\Omega\) 满足
\(
\Omega(\mathsf T_6u,v)=\Omega(u,\mathsf T_6v)
\)，
则对任意特征向量 \(u\in E_\lambda\)、\(v\in E_\mu\) 有
\[
\lambda\,\Omega(u,v)=\Omega(\mathsf T_6u,v)=\Omega(u,\mathsf T_6v)=\mu\,\Omega(u,v).
\]
当 \(\lambda\neq\mu\) 时，推出 \(\Omega(u,v)=0\)。
当 \(\lambda=\mu\) 时，由交替性 \(\Omega(u,u)=0\)，且 \(E_\lambda\) 一维，故同样有 \(\Omega|_{E_\lambda\times E_\lambda}=0\)。
于是 \(\Omega\) 在全部特征基上恒为零，不可能非退化。第 \(2\) 条成立。证毕。
\end{proof}

\begin{theorem}[最小一阶 Markov 障碍纯属 \texorpdfstring{$F_8$}{F8} 尾位通道]\label{thm:window6-minimal-markov-obstruction-pure-f8-tail}
\leanverified{paper\_window6\_minimal\_markov\_obstruction\_pure\_f8\_tail}
记
\[
b^{\mathrm{bin}}_6(\omega):=z_7\in\{0,1\},
\]
即 \(\omega\in\Omega_6\) 对应 Zeckendorf 数字中的首个被截断尾位（权重 \(F_8=21\)）。
则：
\begin{enumerate}
  \item 同纤维两点
  \[
  \omega_0=\texttt{000000},\qquad
  \omega_1=\texttt{010101}
  \]
  满足
  \[
  F_6(\omega_0)=F_6(\omega_1)=\texttt{000000},
  \qquad
  b^{\mathrm{bin}}_6(\omega_0)=0,\quad
  b^{\mathrm{bin}}_6(\omega_1)=1.
  \]
  \item 这对同纤维点的一步邻域在稳定类型上的计数分布不同，因此任何强可并细化都必须将其分离。
  \item 二者之差恰为
  \[
  \omega_0\oplus\omega_1=\texttt{010101},
  \qquad
  \mathrm{int}_6(\texttt{010101})=21=F_8.
  \]
  \item 另一方面，window-\(6\) 的全部几何二点方向仅为
  \[
  \texttt{100010},\qquad \texttt{100110},\qquad \texttt{111110},
  \]
  从而 \(\texttt{010101}\) 不属于任何几何二点方向谱。
\end{enumerate}
因此，window-\(6\) 的最小一阶 Markov 修复完全落在被截断的 \(F_8\) 尾位通道；它既不是 boundary 的 \(F_9\)-sheet parity，也不是任何几何稳定子可见的二点换位方向。
\end{theorem}

\begin{proof}
命题~\ref{prop:terminal-foldbin6-strong-lumpability-fails} 与推论~\ref{cor:terminal-foldbin6-strong-lumpability-binary-witness} 已给出：\(\omega_0=\texttt{000000}\) 与 \(\omega_1=\texttt{010101}\) 落在同一 visible fiber，却具有不同的一步邻域计数，而二值见证恰为
\(
b^{\mathrm{bin}}_6=z_7
\)
；这证明了前两条。

又由
\(
\texttt{000000}\oplus\texttt{010101}=\texttt{010101}
\)
且
\(
\mathrm{int}_6(\texttt{010101})=21=F_8
\)，
得到第 \(3\) 条。
最后，定理~\ref{thm:terminal-foldbin6-two-point-fiber-direction-spectrum} 证明 window-\(6\) 的全部几何二点方向仅为
\(\texttt{100010}\)、\(\texttt{100110}\)、\(\texttt{111110}\)，
因此 \(\texttt{010101}\) 不属于几何方向谱。证毕。
\end{proof}

\begin{theorem}[\texorpdfstring{$F_8$}{F8} 修复位不存在 boundary-covariant 延拓]\label{thm:window6-f8-repair-bit-no-boundary-covariant-extension}
\leanverified{paper\_window6\_f8\_repair\_bit\_no\_boundary\_covariant\_extension}
记
\[
B:=\bigcup_{w\in X_6^{\mathrm{bdry}}}F_6^{-1}(w)
\]
为 boundary 三条二点纤维之并。
称 bit-valued 映射 \(\chi:B\to\{0,1\}\) 为 boundary-covariant，若对每个 \(w\in X_6^{\mathrm{bdry}}\) 与其纤维内换位 \(\tau_w\) 都有
\[
\chi(\tau_w\eta)=1-\chi(\eta)\qquad(\eta\in F_6^{-1}(w)).
\]
则 \(b^{\mathrm{bin}}_6|_B\) 不是 boundary-covariant；更强地，\(b^{\mathrm{bin}}_6\) 在每条 boundary 纤维上都恒为 \(0\)。
\end{theorem}

\begin{proof}
由推论~\ref{cor:terminal-foldbin6-boundary-pure-f9-alias}，对任意 \(w\in X_6^{\mathrm{bdry}}\) 都有
\[
F_6^{-1}(w)=\{V_6(w),\,V_6(w)+34\},
\]
即 boundary 纤维的非平凡尾部别名纯粹发生在 \(F_9=34\) 通道。
再由命题~\ref{prop:terminal-foldbin6-tail-cube-section} 的 tail-cube 描述，boundary 情形满足 \(w_6=1\)，因而 \(F_8\)-尾位对应的变量 \(c_7\) 被强制为 \(0\)。
于是
\[
b^{\mathrm{bin}}_6=z_7=0
\]
在每条 boundary 纤维上恒定。
这显然不可能满足
\(
\chi(\tau_w\eta)=1-\chi(\eta)
\)
的 boundary-covariant 反对称关系。证毕。
\end{proof}

\begin{corollary}[\texorpdfstring{$F_8$}{F8} 修复位与 canonical boundary parity 属于不同隐藏通道]\label{cor:window6-f8-repair-vs-f9-boundary-parity-separated}
\leanverified{paper\_window6\_f8\_repair\_vs\_f9\_boundary\_parity\_separated}
首个被截断的 \(F_8\) 修复位 \(b^{\mathrm{bin}}_6\) 与 canonical boundary sheet parity 不是同一个隐藏量的两种表述。
前者没有 boundary-covariant 延拓，后者则恰由 boundary 纤维上的反对称换位结构定义。
因此二者在 window-\(6\) 中属于不同的隐藏通道：\(F_8\) 通道只负责离散层的一阶 Markov 闭合，而 \(F_9\) 通道则编码 boundary sheet-flip 的规范 \(\ZZ_2^3\) 结构。
\end{corollary}

\begin{proof}
定理~\ref{thm:window6-f8-repair-bit-no-boundary-covariant-extension} 已表明 \(b^{\mathrm{bin}}_6|_B\) 在 boundary 纤维上恒定，故不能承担任何反对称 sheet-flip 角色。
另一方面，推论~\ref{cor:fold-bin-boundary-center-z2} 与推论~\ref{cor:terminal-delta-parity-rule} 已给出：boundary 扇区的 canonical parity 正是由 \(F_9=34\) 二层别名上的纤维内换位反对称结构所定义。
故二者不能互相替代。证毕。
\end{proof}

\begin{conjecture}[Archimedean 选模与 \texorpdfstring{$F_8/F_9$}{F8/F9} 双隐通道分离]\label{conj:window6-archimedean-mode-selection-f8-f9-separation}
设某一临界六顶点系综的 window-\(6\) 可见观测量在适当归一化下出现非平凡 Gaussian free field 极限。
则与 window-\(6\) 相容的连续机制应满足下列四条：
\begin{enumerate}
  \item 可见 massless mode 的选取必经由 Archimedean 频带切片，而不能由任何 \(\QQ\)-有理谱投影实现；
  \item 该 massless mode 不携带由 \(\mathsf T_6\) 直接遗传而来的复极化或辛极化，因此主导连续场首先表现为实标量；
  \item \(F_8\) 修复位只负责离散层的一阶 Markov 闭合，并在 GFF 标度下洗出为局域接触修正；
  \item \(F_9\)-boundary parity 则保留为紧致的 \(\ZZ_2^3\) 拓扑缺陷，而不提升为第二个连续玻色子。
\end{enumerate}
\end{conjecture}
